Guo et al.~\cite{guo2017calibration} showed that modern networks are over-confident and introduced temperature scaling (TS), a one-parameter post-hoc rescaling of the logits fitted on held-out in-distribution data, together with the ECE/NLL evaluation protocol we use. Ovadia et al.~\cite{ovadia2019can} then asked how predictive uncertainty behaves under dataset shift, comparing a plain network, TS, MC dropout, deep ensembles and other methods over graded shift intensities on image corruptions and other data; their headline finding was that quality degrades with shift and that ensembles were the most robust of the methods they compared. Our setup is a small-scale analogue of that protocol, with corruption-style shift intensities in the spirit of~\cite{hendrycks2019benchmarking}. Tomani et al.~\cite{tomani2020post} argue that post-hoc calibration fitted in-domain breaks under domain drift and propose to fit it on perturbed calibration data; our oracle-temperature ablation is a cheap version of that idea, using validation data transformed exactly like the test condition.

Deep ensembles~\cite{lakshminarayanan2016simple} average independently trained networks; MC dropout~\cite{gal2015dropout} keeps dropout active at test time and averages stochastic passes as an approximation to Bayesian inference, a reading that Le~Folgoc et al.~\cite{folgoc2021is} question. Uncertainty Baselines~\cite{nado2021uncertainty} provides tuned large-scale implementations of these and related methods; Laplace approximations~\cite{daxberger2021laplace}, SWAG~\cite{maddox2019simple} and architecture-diverse ensembles~\cite{zaidi2020neural} are further approaches that we do not evaluate. We do not propose a new method; the contribution is a controlled, fully reproducible CPU-scale comparison with a mechanism ablation, and its conclusions should be read at that scale.
